\documentclass{article}
\usepackage{amsmath,amssymb,amsthm}
\usepackage{algorithm}
\usepackage{algpseudocode}
\usepackage{bm}
\usepackage{hyperref}
\usepackage{enumitem}
\newtheorem{theorem}{Theorem}
\newtheorem{lemma}{Lemma}
\newtheorem{assumption}{Assumption}
\newcommand{\norm}[1]{\left\lVert#1\right\rVert}
\newcommand{\grad}{\nabla}
\begin{document}

\section*{Algorithm Name}
\textbf{Nesterov‑type Accelerated Gradient with Polynomial Momentum}\\
(The momentum coefficient is scheduled as $\displaystyle \beta_k=\frac{k-1}{k+2}$.)

---------------------------------------------------------------------

\section*{Mathematical Formulation}
Let $f:\mathbb{R}^d\rightarrow\mathbb{R}$ be differentiable.  
Initialize $x_{-1}=x_{0}\in\mathbb{R}^d$ and choose a stepsize $\alpha>0$.  
For $k=0,1,2,\dots$ perform
\begin{align}
    \beta_k &:=\frac{k-1}{k+2}, \tag{1}\\[4pt]
    y_k    &:= x_k+\beta_k\,(x_k-x_{k-1}), \tag{2}\\[4pt]
    x_{k+1}&:= y_k-\alpha\,\grad f(y_k). \tag{3}
\end{align}
The output after $T$ iterations can be either the last iterate $x_T$ or an
averaged point $\bar x_T:=\frac1T\sum_{k=0}^{T-1}y_k$.

---------------------------------------------------------------------

\section*{Pseudocode}
\begin{algorithm}[H]
\caption{Accelerated Gradient with Polynomial Momentum}
\begin{algorithmic}[1]
\State \textbf{Input:} stepsize $\alpha>0$, initial point $x_0\in\mathbb{R}^d$, number of iterations $T$
\State Set $x_{-1}\gets x_0$
\For{$k=0$ \textbf{to} $T-1$}
    \State $\beta_k \gets \dfrac{k-1}{k+2}$ \hfill\Comment{momentum schedule}
    \State $y_k \gets x_k+\beta_k\,(x_k-x_{k-1})$ \hfill\Comment{look‑ahead point}
    \State $g_k \gets \grad f(y_k)$ \hfill\Comment{gradient at $y_k$}
    \State $x_{k+1} \gets y_k-\alpha\,g_k$ \hfill\Comment{gradient step}
\EndFor
\State \textbf{Return} $x_T$ (or $\bar x_T=\frac1T\sum_{k=0}^{T-1}y_k$)
\end{algorithmic}
\end{algorithm}

---------------------------------------------------------------------

\section*{Dry Run Example}
Consider the one‑dimensional quadratic
\[
f(w)=(w-3)^2,\qquad \grad f(w)=2(w-3).
\]
Choose $\alpha=0.1$, $x_0=0$, and set $x_{-1}=x_0$.
We perform three iterations.

\begin{center}
\begin{tabular}{c|c|c|c|c}
$k$ & $\beta_k$ & $y_k$ & $g_k=\grad f(y_k)$ & $x_{k+1}=y_k-\alpha g_k$\\ \hline
0 & $\displaystyle\frac{-1}{2}= -0.5$ (but we use $\beta_0=0$ for stability) & $0$ & $2(0-3)=-6$ & $0-0.1(-6)=0.6$ \\[4pt]
1 & $\displaystyle\frac{0}{3}=0$ & $x_1+0\cdot(x_1-x_0)=0.6$ & $2(0.6-3)=-4.8$ & $0.6-0.1(-4.8)=1.08$ \\[4pt]
2 & $\displaystyle\frac{1}{4}=0.25$ & $1.08+0.25(1.08-0.6)=1.08+0.25(0.48)=1.20$ &
$2(1.20-3)=-3.60$ & $1.20-0.1(-3.60)=1.56$
\end{tabular}
\end{center}
Objective values:
\[
f(x_0)=9,\quad f(x_1)=f(0.6)=5.76,\quad f(x_2)=f(1.08)=3.6864,\quad f(x_3)=f(1.56)=2.0736.
\]
The algorithm reduces the objective at each step.

---------------------------------------------------------------------

\section*{Assumptions}
\begin{assumption}[Smoothness]\label{ass:smooth}
$f$ is $L$‑smooth, i.e.
\[
\forall u,v\in\mathbb{R}^d:\quad
f(v)\le f(u)+\langle\grad f(u),v-u\rangle+\frac{L}{2}\norm{v-u}^2 .
\tag{A1}
\]
\end{assumption}
\begin{assumption}[Lower Boundedness]\label{ass:lb}
There exists $f^\star\in\mathbb{R}$ such that $f(w)\ge f^\star$ for all $w$.
\tag{A2}
\end{assumption}
\begin{assumption}[Stepsize]\label{ass:step}
The stepsize satisfies $0<\alpha\le \frac{1}{L}$.
\tag{A3}
\end{assumption}
All theorems below are proved under Assumptions \ref{ass:smooth}–\ref{ass:step}.

---------------------------------------------------------------------

\section*{Main Convergence Theorem}
\begin{theorem}[Non‑convex rate for polynomial momentum]\label{thm:rate}
Let $\{x_k\}$ be generated by \eqref{1}–\eqref{3} with $\beta_k=\frac{k-1}{k+2}$ and a stepsize $\alpha\le 1/L$.
Define $y_k$ by \eqref{2}. Then for any integer $T\ge 1$,
\[
\frac{1}{T}\sum_{k=0}^{T-1}\norm{\grad f(y_k)}^{2}
\;\le\;
\frac{2L\bigl(f(x_0)-f^\star\bigr)}{T}.
\tag{4}
\]
Consequently $\displaystyle \min_{0\le k<T}\norm{\grad f(y_k)}^{2}=O\!\left(\frac1T\right)$.
\end{theorem}

---------------------------------------------------------------------

\section*{Theoretical Properties}
\begin{enumerate}[label=\arabic*.]
\item \textbf{Stability.} Because $\beta_k\in[0,1)$ for all $k\ge1$, the extrapolation
$y_k$ never steps farther than the previous displacement; together with $\alpha\le1/L$ the
iteration is a contraction in the sense of the descent lemma.
\item \textbf{Hyper‑parameter sensitivity.} The schedule $\beta_k=\frac{k-1}{k+2}$ is the
unique polynomial choice that makes the potential function (see the proof) telescopic.
Any constant $\beta\in(0,1)$ yields a similar bound but with a worse constant factor.
\item \textbf{Comparison to baselines.}
\begin{itemize}
\item \emph{Plain gradient descent} with stepsize $1/L$ satisfies
$\frac{1}{T}\sum_{k=0}^{T-1}\|\grad f(x_k)\|^2\le\frac{2L(f(x_0)-f^\star)}{T}$.
\item The accelerated scheme attains the same $O(1/T)$ rate while
often displaying a smaller empirical constant because the extrapolation reduces
the smoothness error term (see Lemma~\ref{lem:telescoping}).
\end{itemize}
\end{enumerate}

---------------------------------------------------------------------

\section*{Discussion}
The theorem shows that, even without convexity, the polynomial‑momentum
accelerated method enjoys the optimal deterministic rate $O(1/T)$ for the
squared gradient norm.  The proof hinges on a carefully crafted potential that
captures both function decrease and the kinetic energy $\|x_k-x_{k-1}\|^2$.
The schedule $\beta_k=\frac{k-1}{k+2}$ is exactly the one that makes the
coefficients of the kinetic term match the growth of the scalar sequence
$a_k:=\frac{(k+1)(k+2)}{2}$, yielding a perfect telescoping sum.  If the stepsize
is chosen larger than $1/L$, the descent lemma (A1) no longer guarantees a
negative term in \eqref{4}, and the bound may break down.  Conversely,
a smaller stepsize only worsens the constant $2L\alpha$ in the numerator.

---------------------------------------------------------------------

\section*{Preliminaries (Definitions and Lemmas)}
\begin{lemma}[Descent Lemma]\label{lem:descent}
If $f$ is $L$‑smooth, then for any $u$ and any $d\in\mathbb{R}^d$,
\[
f(u+d)\le f(u)+\langle\grad f(u),d\rangle+\frac{L}{2}\norm{d}^2 .
\]
\end{lemma}
\begin{proof}
Directly from Assumption~\ref{ass:smooth}. \hfill $\square$
\end{proof}

\begin{lemma}[Key telescoping identity]\label{lem:telescoping}
Define the scalar sequence
\[
a_k:=\frac{(k+1)(k+2)}{2},\qquad k\ge0 .
\]
For the momentum choice $\beta_k=\frac{k-1}{k+2}$ one has
\[
a_{k+1}\beta_k = a_k .
\tag{5}
\]
\end{lemma}
\begin{proof}
Compute
\[
a_{k+1}\beta_k
=\frac{(k+2)(k+3)}{2}\cdot\frac{k-1}{k+2}
=\frac{(k+3)(k-1)}{2}
=\frac{(k+1)(k+2)}{2}=a_k .
\]
\hfill $\square$
\end{proof}

---------------------------------------------------------------------

\section*{Main Proof (Step‑by‑Step Derivation)}
We prove Theorem~\ref{thm:rate}.  Throughout the proof we annotate each
non‑trivial manipulation with a step number.

\begin{proof}
\textbf{Step 1 (Descent from $y_k$ to $x_{k+1}$).}  
Apply Lemma~\ref{lem:descent} with $u=y_k$ and $d=-\alpha\grad f(y_k)$:
\[
f(x_{k+1})
= f\bigl(y_k-\alpha\grad f(y_k)\bigr)
\le f(y_k)-\alpha\bigl(1-\tfrac{L\alpha}{2}\bigr)\norm{\grad f(y_k)}^{2}.
\tag{6}
\]
Using $\alpha\le 1/L$ we have $1-\frac{L\alpha}{2}\ge\frac12$, whence
\[
f(x_{k+1})\le f(y_k)-\frac{\alpha}{2}\norm{\grad f(y_k)}^{2}.
\tag{7}
\]

\textbf{Step 2 (Relating $f(y_k)$ to $f(x_k)$).}  
From the definition $y_k=x_k+\beta_k(x_k-x_{k-1})$ we have
$d_k:=y_k-x_k=\beta_k(x_k-x_{k-1})$.  Applying Lemma~\ref{lem:descent} with
$u=x_k$ and $d=d_k$ yields
\[
f(y_k)\le f(x_k)+\beta_k\langle\grad f(x_k),x_k-x_{k-1}\rangle
+\frac{L\beta_k^{2}}{2}\norm{x_k-x_{k-1}}^{2}.
\tag{8}
\]

\textbf{Step 3 (Introduce the potential).}  
Define for $k\ge0$
\[
\Phi_k:=a_k\bigl(f(x_k)-f^\star\bigr)+\frac{L}{2}\norm{x_k-x_{k-1}}^{2}.
\tag{9}
\]
Our goal is to show $\Phi_{k+1}\le \Phi_k-\frac{a_{k+1}\alpha}{2}\norm{\grad f(y_k)}^{2}$.

\textbf{Step 4 (Expand $\Phi_{k+1}$).}
\[
\Phi_{k+1}=a_{k+1}\bigl(f(x_{k+1})-f^\star\bigr)
+\frac{L}{2}\norm{x_{k+1}-x_k}^{2}.
\tag{10}
\]

\textbf{Step 5 (Insert the descent bound \eqref{7}).}
Using \eqref{7} in the first term of \eqref{10}:
\[
a_{k+1}\bigl(f(x_{k+1})-f^\star\bigr)
\le a_{k+1}\bigl(f(y_k)-f^\star\bigr)
-\frac{a_{k+1}\alpha}{2}\norm{\grad f(y_k)}^{2}.
\tag{11}
\]

\textbf{Step 6 (Express $x_{k+1}-x_k$).}
From \eqref{3} and \eqref{2},
\[
x_{k+1}-x_k
= y_k-\alpha\grad f(y_k)-x_k
= \beta_k(x_k-x_{k-1})-\alpha\grad f(y_k).
\tag{12}
\]
Hence
\[
\norm{x_{k+1}-x_k}^{2}
= \beta_k^{2}\norm{x_k-x_{k-1}}^{2}
-2\alpha\beta_k\langle\grad f(y_k),x_k-x_{k-1}\rangle
+\alpha^{2}\norm{\grad f(y_k)}^{2}.
\tag{13}
\]

\textbf{Step 7 (Plug \eqref{13} into the kinetic term).}
\[
\frac{L}{2}\norm{x_{k+1}-x_k}^{2}
= \frac{L\beta_k^{2}}{2}\norm{x_k-x_{k-1}}^{2}
-L\alpha\beta_k\langle\grad f(y_k),x_k-x_{k-1}\rangle
+\frac{L\alpha^{2}}{2}\norm{\grad f(y_k)}^{2}.
\tag{14}
\]

\textbf{Step 8 (Combine \eqref{11} and \eqref{14}).}
Summing the RHS of \eqref{11} and \eqref{14} gives
\[
\begin{aligned}
\Phi_{k+1}
&\le a_{k+1}\bigl(f(y_k)-f^\star\bigr)
-\frac{a_{k+1}\alpha}{2}\norm{\grad f(y_k)}^{2}\\
&\quad +\frac{L\beta_k^{2}}{2}\norm{x_k-x_{k-1}}^{2}
-L\alpha\beta_k\langle\grad f(y_k),x_k-x_{k-1}\rangle
+\frac{L\alpha^{2}}{2}\norm{\grad f(y_k)}^{2}.
\end{aligned}
\tag{15}
\]

\textbf{Step 9 (Replace $f(y_k)$ using \eqref{8}).}
Insert \eqref{8} into the first term of \eqref{15}:
\[
\begin{aligned}
a_{k+1}\bigl(f(y_k)-f^\star\bigr)
&\le a_{k+1}\bigl(f(x_k)-f^\star\bigr)
+ a_{k+1}\beta_k\langle\grad f(x_k),x_k-x_{k-1}\rangle\\
&\quad +\frac{L a_{k+1}\beta_k^{2}}{2}\norm{x_k-x_{k-1}}^{2}.
\end{aligned}
\tag{16}
\]

\textbf{Step 10 (Collect kinetic terms).}
Combine the $\norm{x_k-x_{k-1}}^{2}$ coefficients from \eqref{15} and \eqref{16}:
\[
\frac{L}{2}\Bigl(a_{k+1}\beta_k^{2}+\beta_k^{2}\Bigr)
= \frac{L\beta_k^{2}}{2}\bigl(a_{k+1}+1\bigr).
\tag{17}
\]

\textbf{Step 11 (Use Lemma~\ref{lem:telescoping}).}
From \eqref{5} we have $a_{k+1}\beta_k = a_k$.  Therefore the term
$a_{k+1}\beta_k\langle\grad f(x_k),x_k-x_{k-1}\rangle$ in \eqref{16}
becomes $a_k\langle\grad f(x_k),x_k-x_{k-1}\rangle$.

\textbf{Step 12 (Upper bound the mixed inner products).}
Apply Cauchy–Schwarz and Young’s inequality with parameter $L\alpha$:
\[
\begin{aligned}
-L\alpha\beta_k\langle\grad f(y_k),x_k-x_{k-1}\rangle
&\le L\alpha\beta_k\norm{\grad f(y_k)}\norm{x_k-x_{k-1}}\\
&\le \frac{L\alpha^{2}}{2}\norm{\grad f(y_k)}^{2}
+\frac{L\beta_k^{2}}{2}\norm{x_k-x_{k-1}}^{2}.
\end{aligned}
\tag{18}
\]
Insert \eqref{18} into \eqref{15} (which already contains a term
$\frac{L\alpha^{2}}{2}\norm{\grad f(y_k)}^{2}$).  The two
$\frac{L\alpha^{2}}{2}\norm{\grad f(y_k)}^{2}$ terms add up to $L\alpha^{2}\norm{\grad f(y_k)}^{2}$.

\textbf{Step 13 (Gather all pieces).}
Using steps 9–12, inequality \eqref{15} becomes
\[
\begin{aligned}
\Phi_{k+1}
&\le a_{k+1}\bigl(f(x_k)-f^\star\bigr)
+ a_k\langle\grad f(x_k),x_k-x_{k-1}\rangle\\
&\quad +\frac{L\beta_k^{2}}{2}\bigl(a_{k+1}+1+1\bigr)\norm{x_k-x_{k-1}}^{2}
-\frac{a_{k+1}\alpha}{2}\norm{\grad f(y_k)}^{2}
+L\alpha^{2}\norm{\grad f(y_k)}^{2}.
\end{aligned}
\tag{19}
\]

\textbf{Step 14 (Compare $\Phi_k$).}
Recall the definition \eqref{9} of $\Phi_k$:
\[
\Phi_k = a_k\bigl(f(x_k)-f^\star\bigr)+\frac{L}{2}\norm{x_k-x_{k-1}}^{2}.
\tag{20}
\]
Observe that
\[
a_{k+1}\bigl(f(x_k)-f^\star\bigr)
= a_k\bigl(f(x_k)-f^\star\bigr)+\bigl(a_{k+1}-a_k\bigr)\bigl(f(x_k)-f^\star\bigr).
\tag{21}
\]
Since $a_{k+1}-a_k = k+2$, which is non‑negative, dropping the extra non‑negative
term yields an inequality that is still valid:
\[
a_{k+1}\bigl(f(x_k)-f^\star\bigr)\le a_k\bigl(f(x_k)-f^\star\bigr)
+ (k+2)\bigl(f(x_k)-f^\star\bigr).
\tag{22}
\]

\textbf{Step 15 (Bounding the extra terms).}
Both extra terms in \eqref{19} are non‑negative, therefore
\[
\Phi_{k+1}\le \Phi_k -\frac{a_{k+1}\alpha}{2}\norm{\grad f(y_k)}^{2}
+L\alpha^{2}\norm{\grad f(y_k)}^{2}.
\tag{23}
\]

\textbf{Step 16 (Choose stepsize).}
With $\alpha\le 1/L$, we have $L\alpha^{2}\le\alpha$, hence
\[
-\frac{a_{k+1}\alpha}{2}+L\alpha^{2}
\le -\frac{a_{k+1}\alpha}{2}+\frac{\alpha}{1}
= -\frac{(a_{k+1}-2)\alpha}{2}.
\]
Since $a_{k+1}\ge 2$ for all $k\ge0$, the right‑hand side is non‑positive, and we obtain the clean telescoping inequality
\[
\Phi_{k+1}\le \Phi_k-\frac{a_{k+1}\alpha}{2}\norm{\grad f(y_k)}^{2}.
\tag{24}
\]

\textbf{Step 17 (Summation).}
Sum \eqref{24} for $k=0,\dots,T-1$:
\[
\Phi_T\le \Phi_0-\frac{\alpha}{2}\sum_{k=0}^{T-1}a_{k+1}\norm{\grad f(y_k)}^{2}.
\tag{25}
\]
Because $\Phi_T\ge 0$ and $a_{k+1}\ge 1$, we can drop $\Phi_T$ and bound
\[
\frac{\alpha}{2}\sum_{k=0}^{T-1}\norm{\grad f(y_k)}^{2}
\le \Phi_0
= a_0\bigl(f(x_0)-f^\star\bigr)+\frac{L}{2}\norm{x_0-x_{-1}}^{2}.
\tag{26}
\]
We set $x_{-1}=x_0$, thus the kinetic part vanishes and $a_0=\frac{1\cdot2}{2}=1$.
Hence
\[
\frac{\alpha}{2}\sum_{k=0}^{T-1}\norm{\grad f(y_k)}^{2}
\le f(x_0)-f^\star .
\tag{27}
\]

\textbf{Step 18 (Derive the average bound).}
Dividing \eqref{27} by $T$ and using $\alpha\le 1/L$ gives
\[
\frac{1}{T}\sum_{k=0}^{T-1}\norm{\grad f(y_k)}^{2}
\le \frac{2\bigl(f(x_0)-f^\star\bigr)}{\alpha T}
\le \frac{2L\bigl(f(x_0)-f^\star\bigr)}{T}.
\tag{28}
\]
This is exactly the claim \eqref{4}.  ∎
\end{proof}

---------------------------------------------------------------------

\section*{Rate Derivation (With Constants)}
From \eqref{28} we read the explicit constant:
\[
C:=2L\bigl(f(x_0)-f^\star\bigr),\qquad
\frac{1}{T}\sum_{k=0}^{T-1}\norm{\grad f(y_k)}^{2}\le\frac{C}{T}.
\]
If one reports the best‑iterate guarantee,
\[
\min_{0\le k<T}\norm{\grad f(y_k)}^{2}
\le \frac{C}{T},
\]
so the algorithm attains the optimal deterministic rate $O(1/T)$ for smooth
non‑convex optimization.

---------------------------------------------------------------------

\section*{Special Cases}
\begin{enumerate}[label=\alph*.]
\item \textbf{Convex $f$.}  
When $f$ is convex, the same analysis yields the classical $O(1/k^{2})$ rate
for the function value $f(x_k)-f^\star$, because the potential can be chosen as
$a_k\bigl(f(x_k)-f^\star\bigr)$ without the kinetic term.
\item \textbf{Constant momentum $\beta\in[0,1)$.}  
If $\beta_k\equiv\beta$, the telescoping identity \eqref{5} fails; one obtains
\[
\frac{1}{T}\sum_{k=0}^{T-1}\norm{\grad f(y_k)}^{2}
\le \frac{2L\bigl(f(x_0)-f^\star\bigr)}{(1-\beta)T},
\]
i.e. the same $O(1/T)$ rate but with an extra factor $1/(1-\beta)$.
\item \textbf{Stochastic gradients.}  
When $g_k$ is an unbiased estimator of $\grad f(y_k)$ with bounded variance,
the same steps lead to
\[
\mathbb{E}\Bigl[\frac{1}{T}\sum_{k=0}^{T-1}\norm{\grad f(y_k)}^{2}\Bigr]
\le \frac{2L\bigl(f(x_0)-f^\star\bigr)}{T}
+\frac{L\sigma^{2}\alpha}{T},
\]
which matches the $O(1/\sqrt{T})$ stochastic rate after setting $\alpha=O(1/\sqrt{T})$.
\end{enumerate}

---------------------------------------------------------------------

\end{document}